\sectionkicker{Source-backed technical notes}
\subsection*{Live Webに出さずにAgentを鍛える — World ModelとSynthetic Tools: Technical Notes}
本欄は記事本文で圧縮した一次資料上の情報を、比較・再検証しやすい形で整理したものである。時系列、確認済みの主張、留意点、一次資料URLを掲載する。
\medskip
\subsection*{Theme at a glance}
\addcontentsline{toc}{subsection}{Theme at a glance}
\begin{center}
\footnotesize
\begin{tabularx}{\linewidth}{@{}p{0.20\linewidth}p{0.10\linewidth}p{0.14\linewidth}X@{}}
\toprule
資料 & 位置づけ & 種別 & 時系列 \\
\midrule
DynaWeb: Model-Based Reinforcement Learning of Web Agents & 主要資料 & 論文 & 2026-01-29 (論文投稿) \\
SYNTHAgent: synthetic tasks, simulated tools, rubric rewards & 主要資料 & 論文 & 2026-01-30 (論文投稿) \\
\bottomrule
\end{tabularx}
\end{center}
\normalsize

\begin{technicalnote}{DynaWeb: Model-Based Reinforcement Learning of Web Agents}{主要資料}
\begingroup
% reader-facing Technical Notes break policy
\widowpenalty=10000
\clubpenalty=10000
\displaywidowpenalty=10000
\begin{tabularx}{\linewidth}{@{}>{\bfseries}p{0.22\linewidth}X@{}}
組織 & Hang Ding et al. \\
種別 & 論文 \\
時系列 & 2026-01-29 (論文投稿) \\
\end{tabularx}
\smallskip
{\bfseries 一次資料から整理したtechnical points}
\begin{itemize}[leftmargin=1.5em,itemsep=0.35em]
\item \textbf{一次情報で確認できる事実}: DynaWebはAgent actionからpage representationを予測するweb world modelを学習し、それをmodel-based reinforcement learningのrollout用synthetic environmentとして利用する。
\item \textbf{Author claim}: 論文はopen-source web-agent modelについてWebArenaとWebVoyagerでの改善を報告しているが、これは著者報告の実験結果である。
\end{itemize}
\begin{minipage}{\linewidth}
% reader-facing Technical Notes generic boundary/source tail group
{\bfseries 読む際の境界}
\begin{itemize}[leftmargin=1.5em,itemsep=0.35em]
\item \textbf{分析上の留意点}: simulated web world modelは、敵対的かつ動的なlive webの挙動と大きく異なる可能性がある。
\item \textbf{分析上の留意点}: locatorは1月後に更新されたarXiv v2であるため、chronologyは1月のoriginal submissionを基準とする。
\end{itemize}
\begin{samepage}
{\bfseries 一次資料}
\begingroup\sloppy
\Urlmuskip=0mu plus 2mu\relax
\begin{itemize}[leftmargin=1.5em,itemsep=0.25em]
\item {\scriptsize\url{http://arxiv.org/abs/2601.22149v2}}
\end{itemize}
\endgroup
\end{samepage}
\end{minipage}
\endgroup
\end{technicalnote}
\medskip

\begin{technicalnote}{SYNTHAgent: synthetic tasks, simulated tools, rubric rewards}{主要資料}
\begingroup
% reader-facing Technical Notes break policy
\widowpenalty=10000
\clubpenalty=10000
\displaywidowpenalty=10000
\begin{tabularx}{\linewidth}{@{}>{\bfseries}p{0.22\linewidth}X@{}}
組織 & Yuanjie Lyu et al. \\
種別 & 論文 \\
時系列 & 2026-01-30 (論文投稿) \\
\end{tabularx}
\smallskip
{\bfseries 一次資料から整理したtechnical points}
\begin{itemize}[leftmargin=1.5em,itemsep=0.35em]
\item \textbf{一次情報で確認できる事実}: SYNTHAgentはtool-use taskとmock tool environmentを合成し、不足するprivate information向けにLLM user simulatorを用い、task-level rubric rewardを構成する。
\item \textbf{Author claim}: 論文は14のmath・search・tool-use datasetで改善を報告し、小型modelが大型baselineを上回る例も示すが、いずれも著者報告値である。
\end{itemize}
\begin{minipage}{\linewidth}
% reader-facing Technical Notes generic boundary/source tail group
{\bfseries 読む際の境界}
\begin{itemize}[leftmargin=1.5em,itemsep=0.35em]
\item \textbf{分析上の留意点}: synthetic environmentは安定性を与える一方、実APIの障害や敵対的条件を過小表現する可能性がある。
\item \textbf{分析上の留意点}: locatorは1月後に更新されたv2であるため、chronologyはoriginal submission dateを用いる。
\end{itemize}
\begin{samepage}
{\bfseries 一次資料}
\begingroup\sloppy
\Urlmuskip=0mu plus 2mu\relax
\begin{itemize}[leftmargin=1.5em,itemsep=0.25em]
\item {\scriptsize\url{http://arxiv.org/abs/2601.22511v2}}
\end{itemize}
\endgroup
\end{samepage}
\end{minipage}
\endgroup
\end{technicalnote}
\medskip
